\documentclass[11pt]{article}
\usepackage[a4paper,margin=1in]{geometry}
\usepackage{subfiles}
\usepackage{amsmath,amssymb,bm}
\usepackage{hyperref}
\usepackage{booktabs}
\usepackage{array}
\usepackage{mathtools}

\newcommand{\rs}{r_{\mathrm{s}}}
\newcommand{\Cc}{\mathcal{C}}

\title{Appendix 16 (to main theory): Scalar Sensitivity of Compact
  Bodies in the Bi-Conformal Framework}
\author{}
\date{}

\begin{document}
\let\phi\varphi
\maketitle

\section*{Purpose and scope}
This appendix analyses the scalar sensitivity parameter~$s$
introduced in the main theory (Sec.~9.3) from the coupled
field equations of the bi-conformal action.  The goals are:
(i)~derive the modified Tolman--Oppenheimer--Volkoff (TOV)
equations for a static spherical star in the bi-conformal
metric; (ii)~record the formal definition of~$s$ in terms of
$M_{\mathrm{Komar}}$ and the kinematic motivation that
selects $s = 1/2$ from the bare-mass Komar integrand
$\rho_0\,e^{-\phi}$ (combining the universal conformal
redshift $e^{\phi/2}$ with the proper-volume factor
$e^{-3\phi/2}$) under the approximate rigid shift; (iii)~examine the status of the structural
``shift-symmetry'' argument that was previously invoked to
promote $s_{\mathrm{excess}} = 0$ to an exact, body-universal
statement; and (iv)~confront the resulting picture with
existing pulsar-timing constraints.  No new fields or
parameters are introduced; every equation follows from the
same action used throughout the main theory.

\paragraph{Epistemic status (this appendix).}
The TOV pressure-gradient equation~\eqref{eq:tov_dp} is
strictly shift-invariant under $\phi\to\phi+\delta$
(Sec.~\ref{sec:interior}).  The coupled scalar
equation~\eqref{eq:phi_interior}, by contrast, acquires an
explicit $e^{-\phi}$ prefactor on its matter-source side and
therefore \emph{breaks} strict shift symmetry in the
presence of matter.  As a result, the earlier claim that
``$s_{\mathrm{excess}} = 0$ exactly and non-perturbatively
for every body'' is \emph{not} a theorem of the classical
coupled system.  What survives is: (a) a strictly
shift-invariant hydrostatic sector; (b) a kinematic
argument from the bare-mass Komar integrand
$\rho_0\,e^{-\phi} = \rho_0\,e^{\phi/2}\,e^{-3\phi/2}$
(combining the universal effective mass
$m_{\mathrm{eff}} = m_0\,e^{\phi/2}$ with the proper
spatial-volume factor $e^{-3\phi/2}$) that motivates
$s=1/2$ at leading order under the approximate rigid
shift; (c) a weak-field regime in which shift invariance
of the full coupled system is recovered approximately.  A fully
rigorous proof that $s_{\mathrm{excess}} = 0$ exactly
(rather than to leading order in $|\phi|$) is an
\emph{open task}, deferred to future work.  All claims
below are labelled accordingly.

\section*{Non-negotiable consistency with the main theory}
\begin{itemize}
\item The covariant action is
\begin{equation}\label{eq:action_app16}
S=\frac{1}{16\pi G}\int d^4x\sqrt{-g}\left[R+\frac12\,
  g^{\mu\nu}\partial_\mu\phi\,\partial_\nu\phi\right]
  +S_m[g_{\mu\nu},\psi]\,,
\end{equation}
with matter minimally coupled to the single physical
metric $g_{\mu\nu}$.
\item The bi-conformal ansatz for a static spherical
configuration is
\begin{equation}\label{eq:metric_static_app16}
  ds^2 = -e^{\phi}\,c^2\,dt^2 + e^{-\phi}
  (dr^2 + r^2\,d\Omega^2)\,.
\end{equation}
\item The scalar field equation in the presence of matter is
$\Box\phi = -(8\pi G/c^4)\,T$, where
$T = g^{\mu\nu}T_{\mu\nu}^{(\mathrm{m})}$.
\end{itemize}

%===================================================================
\section{Interior field equations}\label{sec:interior}
%===================================================================

For a static, spherically symmetric perfect fluid with
rest-mass density~$\rho$ and pressure~$p$ (both measured
locally; the corresponding total energy density is
$\varepsilon=\rho c^2+\text{internal energy}$, but in the
non-relativistic-fluid regime used throughout this appendix
the internal-energy contribution is absorbed into the
equation of state), the matter energy-momentum tensor is
$T^{\mu\nu} = (\rho c^2 + p)\,u^\mu u^\nu/c^2 + p\,g^{\mu\nu}$
[equivalently $(\rho + p/c^2)\,u^\mu u^\nu + p\,g^{\mu\nu}$],
and its trace is $T = -\rho\,c^2 + 3p$.

Substituting the bi-conformal
ansatz~(\ref{eq:metric_static_app16}) into the field equations
and restricting to hydrostatic equilibrium yields two
coupled radial ODEs (the hydrostatic
equation~\ref{eq:tov_dp} for $p(r)$ and the scalar
equation~\ref{eq:phi_interior} for $\phi(r)$), closed by an
equation of state $p=p(\rho)$ and supplemented by an
auxiliary mass function
\begin{equation}\label{eq:mass_func}
  m(r) \equiv \int_0^r 4\pi\,\rho(r')\,e^{-3\phi/2}\,
  {r'}^2\,dr'\,,
\end{equation}
which accounts for the proper spatial volume element
$\sqrt{g^{(3)}} = e^{-3\phi/2}\,r^2\sin\theta$.

\subsection{Modified TOV: pressure gradient}

The radial component of hydrostatic equilibrium,
$\nabla_\mu T^{\mu r} = 0$, gives
\begin{equation}\label{eq:tov_dp}
  \frac{dp}{dr} = -(\rho + p/c^2)\,\frac{c^2}{2}\,
  \frac{d\phi}{dr}\,.
\end{equation}
This differs from the standard TOV equation in that the
gravitational ``acceleration'' is $g = (c^2/2)\,d\phi/dr$
rather than the Schwarzschild expression
$Gm/(r^2(1-2Gm/c^2r))$.  In the weak-field limit
$\phi \approx -\rs/r$ and
$g \approx GM/r^2$, recovering Newton.

\textbf{Key structural property:} Eq.~(\ref{eq:tov_dp})
depends on $d\phi/dr$, \emph{not} on $\phi$ itself.
A constant shift $\phi \to \phi + \delta$ leaves the
pressure gradient, and therefore the entire density and
pressure profiles $\rho(r)$, $p(r)$, unchanged.

\subsection{Scalar field equation inside matter}

Direct reduction of the scalar equation
$\Box\phi = -(8\pi G/c^4)\,T$ in the bi-conformal static
metric~(\ref{eq:metric_static_app16}) gives the
d'Alembertian $\Box\phi = e^{\phi}(\phi''+2\phi'/r)$ and,
with $T = -\rho c^2 + 3p$, yields
\begin{equation}\label{eq:phi_interior}
  \frac{d^2\phi}{dr^2} + \frac{2}{r}\frac{d\phi}{dr}
  = \frac{8\pi G}{c^4}\,e^{-\phi}\,
  (\rho\,c^2 - 3p)\,.
\end{equation}
This form is consistent with the canonical vacuum solution
$\phi = -\rs/r$: at $T=0$ both sides vanish, and
$\phi''+2\phi'/r = 0$ is satisfied identically.  A
symbolic verification of $\Box\phi = 0$ for the vacuum
solution is provided by
\texttt{MAIN/Python/Shapiro/check\_kappa\_consistency.py} (see Issue~\#14
audit).

\textbf{Structural property and partial shift invariance.}
Under a global shift $\phi\to\phi+\delta$ with constant
$\delta$, the left-hand side of
Eq.~\eqref{eq:phi_interior} is invariant (it contains only
derivatives), but the right-hand side acquires a factor
$e^{-\delta}$ through the explicit $e^{-\phi}$ coupling.
The interior scalar equation is therefore \emph{not}
strictly shift-invariant when matter is present, even
for finite constant $\delta$: the multiplicative factor
$e^{-\delta}$ is non-vanishing for any finite shift.  The
shift symmetry is recovered only in the perturbative
small-field limit $|\phi|\ll 1$ (where $e^{-\phi}\approx 1$
and the residual is $O(|\phi|)$ in the sensitivity
expansion); at strong field (neutron-star interiors,
$|\phi|\sim 0.1$ near the centre) the correction is
$\sim 10\%$.

\subsection{Boundary conditions}

At the centre ($r = 0$): $\phi(0) = \phi_c$ (a free
parameter determined by matching), $d\phi/dr|_0 = 0$
(regularity).  At the stellar surface ($r = R$, where
$p = 0$): the interior solution joins the vacuum exterior
$\phi_{\mathrm{ext}} = -\rs/r + \phi_\infty$, where
$\phi_\infty$ is the asymptotic (background) scalar field
value.

%===================================================================
\section{Shift invariance: sectoral analysis}
\label{sec:shift}
%===================================================================

The TOV pressure-gradient
equation~(\ref{eq:tov_dp}) depends only on a
\emph{derivative} of~$\phi$, whereas the reduced scalar
equation~(\ref{eq:phi_interior}) carries an explicit
$e^{-\phi}$ prefactor on its matter-source side.  The two
sectors therefore have different behaviour under a rigid
shift $\phi\to\phi+\delta$, and must be analysed
separately.

\subsection{TOV-level shift invariance (rigorous)}
\label{sec:shift_tov}

\begin{quote}
\textbf{Lemma (Hydrostatic shift invariance, exact).}
Fix the scalar profile $\phi_0(r)$ as an externally
specified background (i.e.\ \emph{do not} re-solve the
coupled scalar equation~(\ref{eq:phi_interior}) for the
shifted profile), the equation of state $p=p(\rho)$, and
the central boundary data $\rho(0)$, $p(0)$.  Then the
hydrostatic equation~(\ref{eq:tov_dp}) is invariant
under $\phi_0 \to \phi_0 + \delta$ for any constant
$\delta$: the pressure gradient $dp/dr$, and hence
the density and pressure profiles $\rho(r)$ and $p(r)$
generated by integrating~(\ref{eq:tov_dp}) from the
centre, are unchanged.  The lemma applies to the
hydrostatic sub-equation in isolation; it does
\emph{not} establish invariance of the full coupled
matter--scalar system, which is broken by the
$e^{-\phi}$ source factor of
Eq.~(\ref{eq:phi_interior}).
\end{quote}

\textit{Proof.}---Under $\phi_0 \to \phi_0 + \delta$,
$d\phi_0/dr \to d\phi_0/dr$ (unchanged), so the
right-hand side of~(\ref{eq:tov_dp}) is unchanged.
Integrating from the centre with boundary conditions
at $r=0$, the profiles $\rho(r)$ and $p(r)$ are
identical.~\hfill$\square$

This sub-sectoral invariance is \emph{exact} and
\emph{non-perturbative}.  It is a direct consequence
of the minimal matter coupling to $g_{\mu\nu}$ and of
the constancy of $G_4$ (no Brans--Dicke running of the
gravitational coupling).  It underpins the statement
that ``the stellar structure does not care about the
absolute value of the background scalar, only about
its gradient.''

\subsection{Coupled-system shift symmetry (approximate;
broken at $O(|\phi|)$)}
\label{sec:shift_coupled}

The coupled system is not strictly shift-invariant.
Under $\phi \to \phi + \delta$ with $\rho, p$ held
fixed, the left-hand side of~(\ref{eq:phi_interior})
is unchanged but the right-hand side acquires a factor
$e^{-\delta}$.  To restore a solution, one must rescale
the matter source by $e^{+\delta}$, which is not a
symmetry of the matter action.  Explicitly:
\begin{itemize}
\item in the weak-field limit $|\phi|\ll 1$, one has
  $e^{-\phi} \approx 1 - \phi + \tfrac12\phi^2 - \cdots$,
  and shift invariance is recovered at zeroth order;
\item the leading correction is linear in $|\phi|$;
  for neutron-star interiors $|\phi_c| \sim \rs/R \sim
  0.2$, the shift asymmetry is at the $\sim\!10\%$ level;
\item for vacuum ($\rho = p = 0$), the right-hand side
  vanishes identically and shift invariance is exactly
  restored---consistent with the vacuum no-hair
  derivation of main theory Sec.~7.
\end{itemize}

\textbf{Status.}  An \emph{exact} rigid-shift symmetry
of the coupled TOV--scalar system is not a property of
the classical equations in matter; it holds strictly only
at the level of~(\ref{eq:tov_dp}) and in vacuum.  Any
claim of exact non-perturbative universality of $s$ at
arbitrary compactness must be supported by an argument
beyond a naive rigid shift (e.g.\ an on-shell
reformulation of the matter sector that absorbs the
$e^{-\phi}$ factor into a redefined proper density;
this is the programme flagged in the quantum paper and
discussed in Sec.~\ref{sec:summary_app16} below).
This programme is an \textbf{open task}.

%===================================================================
\section{Scalar sensitivity: formal definition, motivated value}
\label{sec:s_def}
%===================================================================

Place the star in a weak, uniform external scalar field
$\phi_\infty \neq 0$.  The total (Komar) mass measured at
infinity depends on $\phi_\infty$ through two channels:
(a)~the conformal factor $e^{\phi/2}$ that defines the
effective mass (main theory Eq.~(46)); and (b)~the internal
structure adjustment---the star re-equilibrates when the
external field changes.

The scalar sensitivity is
\begin{equation}\label{eq:s_formal}
  \boxed{s \equiv -\frac{1}{2}\,
  \frac{\partial\ln M_{\mathrm{Komar}}}
       {\partial\phi_\infty}
  \bigg|_{\phi_\infty = 0}\,.}
\end{equation}

\subsection{Kinematic motivation: $s = 1/2$ from the
universal conformal rescaling}
\label{sec:s_kinematic}

The position-dependent effective mass assigned to any
localised body in the bi-conformal metric is
\begin{equation}\label{eq:meff_app16}
  m_{\mathrm{eff}}(\phi) = m_0\,e^{\phi/2}\,,
\end{equation}
as derived in the main theory (Sec.~9.2, Eq.~(46)) from
the conformal redshift of the static clock factor.  This
rescaling is \emph{universal}: it multiplies every
localised energy by the same factor $e^{\phi/2}$,
independently of composition or compactness.

The formal sensitivity~(\ref{eq:s_formal}) is defined,
however, in terms of the Komar mass $M_{\mathrm{Komar}}$
that controls the asymptotic-field response, not in terms
of the test-particle Killing-energy
$m_{\mathrm{eff}}$~(\ref{eq:meff_app16}).  Combining the
universal conformal redshift $e^{\phi/2}$ with the proper
spatial-volume element $\sqrt{g^{(3)}}=e^{-3\phi/2}r^2
\sin\theta$ yields the bare-mass Komar
integrand~$\rho_0\,e^{-\phi}$ of
Eq.~\eqref{eq:komar_app16}: the universal redshift and
proper-volume factors combine into the single weight
$e^{\phi/2}\cdot e^{-3\phi/2}=e^{-\phi}$.  If, in the
weak-field/leading-order regime $|\phi|\ll1$, the interior
profile is approximated as shifting rigidly under
$\phi_\infty\to\phi_\infty+\delta$ (an approximation
exact to $O(|\phi|)$, and broken at higher order by the
$e^{-\phi}$ matter source factor of
Eq.~\ref{eq:phi_interior}; see
Sec.~\ref{sec:shift_coupled}), the bare Komar mass
rescales rigidly as
$M_{\mathrm{Komar}}^{(0)}(\phi_\infty+\delta)=
e^{-\delta}\,M_{\mathrm{Komar}}^{(0)}(\phi_\infty)$, so
$\partial\ln M_{\mathrm{Komar}}/\partial\phi_\infty=-1$ at
leading order, and the formal definition~(\ref{eq:s_formal})
gives
\begin{equation}\label{eq:s_kinematic}
  s_{\mathrm{kin}}
  \;=\;-\frac{1}{2}\cdot
  \left.\frac{\partial\ln M_{\mathrm{Komar}}^{(0)}}
       {\partial\phi_\infty}\right|_{\mathrm{rigid\;shift}}
  \;=\;-\frac{1}{2}\cdot(-1)\;=\;\frac{1}{2}\,.
\end{equation}
The kinematic channel therefore selects $s = 1/2$ for
every body at leading order.  The universal
conformal-redshift factor $e^{\phi/2}$ of
Eq.~\eqref{eq:meff_app16} appears here as one of the two
factors in the combined Komar weight $e^{-\phi}$ (alongside
the proper-volume factor $e^{-3\phi/2}$); on its own it
does \emph{not} produce $s=1/2$, and the previous-draft
attempt to differentiate $m_{\mathrm{eff}}$ alone (instead
of $M_{\mathrm{Komar}}$) was a category error in the
formal definition.

\paragraph{Status.}  Equation~(\ref{eq:s_kinematic})
identifies $s=1/2$ as the natural value common to all
bodies from a purely kinematic (geometric) argument.  It
does \emph{not} prove $s_{\mathrm{excess}} = 0$
\emph{exactly}: the structural-adjustment channel (the
star re-equilibrating when $\phi_\infty$ shifts) is not
shown to vanish strictly, because the full coupled shift
symmetry is broken (Sec.~\ref{sec:shift_coupled}).  What
is established here is:
\begin{itemize}
\item in the weak-field regime $|\phi| \ll 1$, the
  structural-response correction is $O(|\phi|)$ and
  sub-dominant, so $s \simeq 1/2 + O(|\phi|)$ for every
  body;
\item at leading order in $|\phi|$, the body-to-body
  difference is $O(|\phi_A - \phi_B|)$ at matched
  asymptotic $\phi_\infty$, vanishing identically when
  the bodies share a common background;
\item the question whether the $O(|\phi|)$ correction
  itself cancels body-to-body at all orders
  (giving $s_{\mathrm{excess}} = 0$ exactly at arbitrary
  compactness) is \textbf{an open task}.
\end{itemize}

\subsection{Structural-response correction (not derived
to vanish)}

The \emph{bare-mass} Komar integral (main theory
Eq.~(47); see also ``Scope of Eq.~(\ref{eq:komar_app16})''
below) is
\begin{equation}\label{eq:komar_app16}
  M_{\mathrm{Komar}}^{(0)} \equiv
  \int \rho_0\,e^{-\phi}\,d^3x\,.
\end{equation}
When $\phi_\infty$ shifts by $\delta$, the interior
profile does \emph{not} shift rigidly: because the
coupled equation~(\ref{eq:phi_interior}) carries the
matter-side factor $e^{-\phi}$, the solution
$\phi(r;\phi_\infty+\delta)$ differs from
$\phi(r;\phi_\infty) + \delta$ by a correction of order
$|\phi|$ localised inside the body.  Consequently the
identity $M_{\mathrm{Komar}}^{(0)}(\phi_\infty+\delta) =
e^{-\delta}\,M_{\mathrm{Komar}}^{(0)}(\phi_\infty)$, which
would give $\partial\ln M_{\mathrm{Komar}}^{(0)}/\partial
\phi_\infty = -1$ exactly, holds only to leading order
in $|\phi|$.

\paragraph{Scope of Eq.~(\ref{eq:komar_app16}):
bare-mass vs.\ Tolman--Komar form.}
Equation~(\ref{eq:komar_app16}) weights the rest-mass
density $\rho_0$ by the gravitational redshift / proper
volume factor $e^{-\phi} = e^{\phi/2}\,e^{-3\phi/2}$
alone.  The stationary Tolman--Komar perfect-fluid mass,
on the other hand, carries an additional pressure
contribution:
\begin{equation}\label{eq:komar_TK}
  M_{\mathrm{Komar}}^{\mathrm{TK}} =
  \int (\rho_0 + 3p/c^2)\,e^{-\phi}\,d^3x\,.
\end{equation}
The two integrals agree in the dust limit $p\ll\rho_0
c^2$ and differ by $O(p/\rho_0 c^2)$ otherwise.  For
neutron-star interiors $p/(\rho_0 c^2)\sim 0.2$--$0.3$
near the centre, the relative difference is $\sim 30\%$
and concentrated in the pressure-dominated core.  The
kinematic motivation of $s=1/2$ (Sec.~\ref{sec:s_kinematic})
uses the full bare-mass weight $e^{-\phi}=e^{\phi/2}\,e^{-3\phi/2}$
(universal conformal redshift combined with proper
spatial-volume factor); both factors are common to the
bare-mass and Tolman--Komar forms, so the motivation is
not altered by this choice.  The body-to-body cancellation of the
residual structural-response correction at
$O(|\phi|)$---the open task flagged in
Sec.~\ref{sec:shift_coupled}---depends on which Komar
integral is adopted for matter-filled strong-field
regimes; the perfect-fluid analysis with
Eq.~(\ref{eq:komar_TK}) is part of the same open task.

\textbf{Working assumption used elsewhere.}  Throughout
the main theory and companion papers, the working value
adopted is
\begin{equation}\label{eq:s_working}
  s = \frac{1}{2}\,,\qquad
  s_{\mathrm{excess}} = 0 \quad
  \text{(leading order; exact result deferred)}.
\end{equation}
The dipole radiation formula (main theory Eq.~(17)) and
the Nordtvedt prediction (main theory Eq.~(26)) are
evaluated with~(\ref{eq:s_working}); the resulting
predictions are correct to leading order in $|\phi|$.
Whether these predictions remain exact at arbitrary
compactness requires the open-task analysis flagged in
Sec.~\ref{sec:shift_coupled}.

\subsection{Physical interpretation (weak-field)}

The approximate vanishing of $s_{\mathrm{excess}}$ has a
transparent physical origin.  Under a rigid background
shift $\phi_\infty\to\phi_\infty+\delta$, the bare-mass
Komar weight $e^{-\phi}=e^{\phi/2}\,e^{-3\phi/2}$
multiplies the Komar integral of every body by the same
overall factor $e^{-\delta}$ at leading order in $|\phi|$
(both the universal conformal-redshift factor $e^{\phi/2}$
and the universal proper-volume factor $e^{-3\phi/2}$
contribute, and the position-dependent test-particle mass
$m_{\mathrm{eff}}=m_0\,e^{\phi/2}$ of main-theory Eq.~(46)
is one of these two factors, not the whole story).  Since
sensitivity measures the \emph{differential} response
between bodies, and the leading-order Komar response
$\partial\ln M_{\mathrm{Komar}}/\partial\phi_\infty=-1$
is the same for every body, the leading-order excess
vanishes.  The higher-order structural correction,
induced by the matter-side $e^{-\phi}$ coupling on the
RHS of Eq.~(\ref{eq:phi_interior}), is not proven to
cancel body-to-body and is the content of the open task.

%===================================================================
\section{Coupled linear perturbation: status}
\label{sec:correction}
%===================================================================

A naive linearisation of the scalar field
equation~(\ref{eq:phi_interior}) around a background
$\phi_0(r)$, \emph{holding the density $\rho$ fixed},
yields an apparent screening equation:
\begin{equation}\label{eq:wrong_pert}
  \frac{1}{r^2}\frac{d}{dr}\!\left(r^2\,
  \frac{d(\delta\phi)}{dr}\right)
  + \frac{8\pi G\rho}{c^2}\,e^{-\phi_0}\,\delta\phi
  = 0\,,
  \qquad\textbf{(incomplete: $\rho$ held fixed)}
\end{equation}
with an effective mass term $\mu^2 \propto 8\pi G\rho\,
e^{-\phi_0}/c^2$ that would screen the external field
perturbation and produce $s_{\mathrm{excess}} \neq 0$.

\textbf{This equation is incomplete.}  It omits the
structural response $(\delta\rho,\,\delta p)$ that
accompanies $\delta\phi$.  The correct coupled
linearisation is recorded below; it weakens but does
\emph{not} fully cancel the screening term at strong
field.

\subsection{Full coupled perturbation}

Let $\phi = \phi_0 + \epsilon\,\delta\phi$,
$\rho = \rho_0 + \epsilon\,\delta\rho$,
$p = p_0 + \epsilon\,\delta p$.  Linearising the
scalar equation~(\ref{eq:phi_interior}) gives:
\begin{align}
  \text{LHS:}\quad &
  \delta\phi'' + \frac{2}{r}\,\delta\phi',
  \label{eq:LHS_pert}\\[4pt]
  \text{RHS:}\quad &
  \frac{8\pi G}{c^4}\,e^{-\phi_0}\,\bigl[
  (\delta\rho\,c^2 - 3\,\delta p)
  - (\rho_0 c^2 - 3p_0)\,\delta\phi\bigr]\,.
  \label{eq:RHS_pert}
\end{align}

\subsection{The rigid-shift ansatz does not solve the
coupled system in matter}

Substitute $\delta\phi = \mathrm{const}$ and
$\delta\rho = \delta p = 0$:
\begin{itemize}
\item LHS: $\delta\phi' = 0 \Rightarrow$ LHS~$=0$.
\item RHS: $\delta\rho = \delta p = 0 \Rightarrow$
  RHS~$= -\delta\phi\,(8\pi G/c^4)\,e^{-\phi_0}\,
  (\rho_0 c^2 - 3p_0)$.
\item The two sides agree only where the background
  matter source $(\rho_0 c^2 - 3p_0)$ vanishes, i.e.\
  outside the star.  Inside the star the rigid shift
  is \emph{not} a solution.
\end{itemize}
This is the linearised counterpart of the strict
shift-symmetry failure in matter established in
Sec.~\ref{sec:shift_coupled}: the naive
rigid-shift ansatz ceases to be an equilibrium once the
explicit $e^{-\phi}$ coupling is retained.

\paragraph{What survives.}  In the vacuum exterior
($\rho_0 = p_0 = 0$), the rigid shift is an exact
solution; in the weak-field regime $|\phi_0| \ll 1$,
the right-hand side is suppressed by $e^{-\phi_0} \approx 1$
and the leading-order rigid-shift picture is restored.
Whether a modified ansatz (e.g.\ one in which the
structural response $(\delta\rho,\delta p)$ is chosen
self-consistently to cancel the matter-side correction)
solves the coupled system non-perturbatively is an open
question: it would amount to showing that the full
coupled system has a hidden rigid-shift invariance not
manifest in the equations as written.  The construction
of such a reformulation (if it exists) is the same
\textbf{open task} referenced in
Sec.~\ref{sec:shift_coupled}.

%===================================================================
\section{No-hair property: scope}
\label{sec:nohair}
%===================================================================

Two statements must be distinguished:
\begin{enumerate}
\item \textbf{Vacuum no-hair (rigorously established).}
  In the vacuum exterior, $\nabla^2\phi = 0$ together with
  asymptotic boundary data fixes the harmonic profile
  $\phi(r) = \phi_\infty - C/r$ uniquely up to the monopole
  coefficient $C$ (main theory Sec.~7 and Appendix~1).
  For an isolated body this profile shifts rigidly under
  $\phi_\infty\to\phi_\infty+\delta$ ($C$ unchanged), so
  the bare-mass Komar
  integral~(\ref{eq:komar_app16}) rescales by the rigid
  factor $e^{-\delta}$ and the formal
  definition~(\ref{eq:s_formal}) gives $s=1/2$,
  $s_{\mathrm{excess}}=0$ as in
  Sec.~\ref{sec:s_kinematic}, but now \emph{exactly}
  rather than at leading order: there is no matter source
  to spoil the rigid shift.  This is fully derived and
  non-perturbative.
\item \textbf{Matter-filled body no-hair (leading order;
  exact statement deferred).}  For stars with internal
  structure, the kinematic argument of
  Sec.~\ref{sec:s_kinematic} selects $s=1/2$ at leading
  order in $|\phi|$.  The full coupled-system shift
  symmetry that would promote this to an exact statement
  at arbitrary compactness is \emph{not} established
  (Sec.~\ref{sec:shift_coupled}); the structural-response
  correction is of order $|\phi_c|$ inside the body and
  is not proven to cancel body-to-body.
\end{enumerate}

\begin{center}
\begin{tabular}{lccl}
\hline
Body type & $s$ & $s_{\mathrm{excess}}$ & Status \\
\hline
Test particle & $1/2$ & $0$ & exact (Komar route, Sec.~\ref{sec:s_kinematic}) \\
Vacuum compact object & $1/2$ & $0$ & exact (rigid-shift Komar response, item~1 above) \\
White dwarf ($\Cc \sim 10^{-4}$) & $\simeq 1/2$ &
  $O(10^{-4})$ & leading order; exact value deferred \\
Neutron star ($\Cc \sim 0.2$) & $\simeq 1/2$ &
  $O(0.2)$ if uncancelled & leading order; exact value deferred \\
\hline
\end{tabular}
\end{center}

\textit{Physical interpretation (weak-field).}---Every
body in the bi-conformal theory responds to a background
scalar field \emph{at leading order} as a test particle
does: its scalar field shifts rigidly with $\phi_\infty$
up to corrections of order $|\phi|$ localised inside the
body.  The leading-order universality reflects that the
internal structure is governed by $d\phi/dr$ (the
gradient) rather than by $\phi$ (the absolute value) at
the level of the TOV pressure-gradient equation.  Whether
the residual $O(|\phi|)$ correction also cancels
body-to-body to yield an exact effacement for strong
field is the open task flagged throughout.

\paragraph{Universal scalar charge, effaced
differential response.}
It is worth emphasising that the no-hair result of this
appendix does \emph{not} say that bodies carry no scalar
charge.  The scalar charge is defined via the response of
the bare-mass Komar integral~(\ref{eq:komar_app16}) to a
rigid shift of the asymptotic scalar value, evaluated at
the reference background $\phi_\infty=0$:
\begin{equation}\label{eq:qs_def}
  q_s \equiv -\frac{\partial M_{\mathrm{Komar}}^{(0)}}
       {\partial\phi_\infty}\bigg|_{\phi_\infty=0}
  = 2\,s\,M\,,
\end{equation}
where $M\equiv M_{\mathrm{Komar}}^{(0)}(\phi_\infty=0)$
is the unshifted Komar mass and the factor $2$ comes
from the convention of~(\ref{eq:s_formal}).  In the
vacuum-exterior monopole profile $\phi(r)=\phi_\infty-C/r$,
the coefficient $C$ is normalised by matching to this
Komar mass: at leading order $C=2GM/c^2=\rs$, so the
boundary scalar charge $C$ and the Komar-response charge
$q_s$ refer to the same quantity in different
parametrisations.  The working value $s=1/2$ (exact for
vacuum compact objects; leading order, all bodies) gives
\begin{equation}\label{eq:qs_universal}
  q_s = M\qquad\text{(leading order, every body)}.
\end{equation}
Every body therefore carries a nonzero scalar charge
proportional to its mass.  What is effaced is not the
charge but the \emph{differential} response: because
$q_s/M$ is a universal constant (independent of
compactness and equation of state at leading order),
binary systems satisfy $s_A = s_B$ and therefore
$(q_{s,A}/M_A) - (q_{s,B}/M_B) = 0$, which kills the
dipole radiation formula of main theory Eq.~(17) and
the Nordtvedt effect.  The distinction
``universal charge with effaced differential response''
is the physically correct statement of the no-hair result;
``no independent scalar charge'' would be a stronger
claim not supported here.

\paragraph{Scope framing: vacuum vs.\ matter derivations.}
The two no-hair statements above are \emph{consistent
across derivations} rather than instances of one proof:
vacuum compact objects follow from the vacuum exterior
harmonic monopole profile
($\phi=\phi_\infty - C/r$, with $C$ fixed by the
mass/source normalization; main theory Sec.~7 and
Appendix~1) \emph{combined with} the
rigid-shift response of the bare-mass Komar
integral~(\ref{eq:komar_app16}); the uniqueness step
alone fixes the profile, the Komar-response step
delivers the sensitivity value (Sec.~\ref{sec:nohair},
item~1).  Matter-filled bodies follow from the kinematic
Komar-integrand argument of Sec.~\ref{sec:s_kinematic}
(rigid-shift response of $\rho_0\,e^{-\phi}=\rho_0\,
e^{\phi/2}\,e^{-3\phi/2}$) at leading order.
Both yield $s=1/2$, $s_{\mathrm{excess}}=0$, but by
different routes; a single unified proof covering both
regimes exactly would require closure of the open task
of Sec.~\ref{sec:shift_coupled}.

%===================================================================
\section{Confrontation with pulsar-timing bounds}
\label{sec:bounds}
%===================================================================

\subsection{Current constraints}

The strongest constraints on scalar sensitivity come from:
\begin{enumerate}
\item \textbf{PSR~J1738+0333} (Freire et~al.\ 2012):
  a pulsar--white-dwarf binary with precisely measured
  orbital decay.  Since the dipole rate scales as
  $(s_{\mathrm{NS}}-s_{\mathrm{WD}})^{2}$ (see
  Eq.~\ref{eq:Pdot_dipole_app16}), the absence of anomalous
  decay constrains the excess sensitivity (relative to the
  universal kinematic value $s=1/2$) to
  $|s_{\mathrm{excess}}|^2 \lesssim 2.5\times10^{-3}$, i.e.\
  $|s_{\mathrm{excess}}| \lesssim 0.05$ at $95\%$
  confidence as an
  \emph{order-of-magnitude illustrative} reading; a
  faithful translation requires the ISPG-specific
  charge-to-coupling normalization for matter-filled
  bodies (open task, see Sec.~\ref{sec:bounds}
  Compatibility-table footnote).

\item \textbf{PSR~J0337+1715} (Archibald et~al.\ 2018):
  a millisecond pulsar in a hierarchical triple with two
  white dwarfs.  Tests the universality of free fall
  (Nordtvedt effect) to
  $|\Delta a/a| < 2.6\times10^{-6}$ (conservative 2018
  Archibald limit; the 2025 Voisin et~al.\
  re-analysis~\cite{Voisin2025} tightens this to
  model-dependent
  $|\Delta|\lesssim 1.5$--$2.3\times10^{-6}$).
\end{enumerate}

\subsection{Compatibility assessment}

Under the working assumption $s_{\mathrm{excess}}\to 0$
(consistent with the leading-order kinematic result of
Sec.~\ref{sec:s_kinematic}), the predictions are:

\begin{center}
\begin{tabular}{lccc}
\hline
Test & Requires$^{\dagger}$ & Predicted & Status \\
\hline
J0337+1715 (Nordtvedt) & $|s_{\mathrm{excess}}| \lesssim 2.6\times10^{-6}$
  & $s_{\mathrm{excess}} \to 0$ & consistent if open task closes \\
J1738+0333 ($\dot{P}_b$) & $|s_{\mathrm{excess}}| \lesssim 0.05$
  & $s_{\mathrm{excess}} \to 0$ & consistent if open task closes \\
LIGO (BH--BH dipole) & $s_{\mathrm{excess}} = 0$
  & $s_{\mathrm{excess}} = 0$ (vacuum exact) & \textbf{passes rigorously} \\
\hline
\end{tabular}
\end{center}
\smallskip\noindent{\footnotesize
$^{\dagger}$The pulsar entries are
\emph{order-of-magnitude illustrative} requirements
obtained by reading published bounds as direct
constraints on $|s_{\mathrm{excess}}|$ in the working
ISPG convention $s\to 1/2+s_{\mathrm{excess}}$.
A faithful comparison to the standard pulsar dipole
coefficient $\kappa_D$ and to the strong-field UFF
parameter $\Delta$ requires deriving the ISPG
charge-to-coupling normalization for matter-filled
bodies; that normalization is part of the same open task
as the matter-filled $s_{\mathrm{excess}}$ cancellation.
The LIGO row is a vacuum-sector statement and is not
subject to that normalization.}

\textit{(a) Orbital decay (J1738+0333).}---The dipole
contribution to the period derivative is (main theory
Eq.~(20)):
\begin{equation}\label{eq:Pdot_dipole_app16}
  \dot{P}_b^{(\mathrm{dipole})} \propto
  (s_{\mathrm{NS}} - s_{\mathrm{WD}})^2\,.
\end{equation}
At leading order in $|\phi|$, $s_{\mathrm{NS}} \simeq
s_{\mathrm{WD}} \simeq 1/2$ and the dipole amplitude is
suppressed.  The residual body-to-body difference is
$O(|\phi_{c,\mathrm{NS}}| - |\phi_{c,\mathrm{WD}}|)
\sim O(0.2)$ if uncancelled, but the observational
bound, read illustratively as a direct constraint on
$|s_{\mathrm{excess}}|$ pending the ISPG-specific
$\kappa_D$ normalization (Sec.~\ref{sec:bounds}
Compatibility-table footnote), allows
$|s_{\mathrm{excess}}| \lesssim 0.05$, so the theory is
consistent with data whenever the cancellation
established at leading order persists at $O(|\phi|)$
after the open task is closed.

\textit{(b) Universality of free fall (J0337+1715).}---The
schematic Nordtvedt-type identity in the ISPG convention
is (main theory Eq.~(26))
\begin{equation}\label{eq:nordtvedt_app16}
  \frac{\Delta a}{a} = s_{\mathrm{NS}} - s_{\mathrm{WD}}\,,
\end{equation}
which reduces, at leading order, to $\Delta a/a \simeq 0$
because $s_{\mathrm{NS}}\simeq s_{\mathrm{WD}}\simeq 1/2$.
The translation of this schematic identity into the
strong-field UFF parameter $\Delta$ measured in pulsar
timing requires an ISPG-specific coupling normalization
for matter-filled bodies (the same open task as the
matter-filled $s_{\mathrm{excess}}$ cancellation); the
identity above is therefore an order-of-magnitude
schematic statement, not a literal substitution into the
published bound.  With that proviso, consistency with
the conservative 2018 Archibald limit
$|\Delta a/a| < 2.6\times10^{-6}$~\cite{Archibald2018}
(later tightened to model-dependent
$|\Delta| \lesssim 1.5$--$2.3\times10^{-6}$ in the 2025
Voisin et al.\ re-analysis~\cite{Voisin2025}) then
requires that the $O(|\phi|)$ structural-response
correction also cancel between NS and WD to
$\sim 10^{-5}$ relative precision.  This stringent bound
is the \emph{empirical} motivation for expecting the
open task to close with $s_{\mathrm{excess}} = 0$
exactly; it is not a derived result of the classical
coupled system.

\textit{(c) LIGO compact-object mergers.}---For vacuum
binaries (BH--BH), the no-hair result holds rigorously:
$\nabla^2\phi = 0$ fixes the exterior harmonic profile
$\phi=\phi_\infty-C/r$ (Sec.~\ref{sec:nohair}, item~1)
and the rigid-shift response of the bare-mass Komar
integral~(\ref{eq:komar_app16}) gives
$s_A = s_B = 1/2$ and $s_{\mathrm{excess}} = 0$ exactly
via Eq.~(\ref{eq:s_formal}).  No dipole radiation is
predicted, consistent with all LIGO--Virgo BH--BH
detections.  For NS--NS and NS--CO binaries, the statement
is leading-order, contingent on the open task.

\subsection{Contrast with Brans--Dicke / DEF}

In Brans--Dicke and DEF-type scalar-tensor theories, the
effective gravitational coupling $G_{\mathrm{eff}}$ depends
on the scalar field through $G_4(\varphi)$.  When the
background field $\phi_\infty$ changes, $G_{\mathrm{eff}}$
changes, forcing the star to restructure.  This gives
$s_{\mathrm{excess}} \neq 0$ with a compactness-dependent
coefficient.

In the present theory, $G_4 = M_{\mathrm{Pl}}^2/2 =
\mathrm{const}$: there is no running of the gravitational
coupling.  The scalar field enters only through the
bi-conformal metric, and the TOV pressure-gradient equation
depends on $d\phi/dr$ rather than on $\phi$.  This
sub-sectoral shift invariance is the structural reason
why $s_{\mathrm{excess}}$ is \emph{small} (leading order
$O(|\phi|)$) rather than $O(1)$, and is what makes the
theory compatible with pulsar-timing data at the
leading-order level.  Whether it is strictly zero at all
orders is the open task.

%===================================================================
\section{Summary}\label{sec:summary_app16}
%===================================================================

\begin{enumerate}
\item The bi-conformal TOV equations~(\ref{eq:tov_dp})--(\ref{eq:phi_interior})
  are the interior counterpart of the exterior exponential
  metric used throughout the main theory.

\item The TOV pressure-gradient equation~(\ref{eq:tov_dp})
  depends only on $d\phi/dr$ and is therefore strictly
  shift-invariant under $\phi\to\phi+\delta$
  (Sec.~\ref{sec:shift_tov}).  The reduced scalar
  equation~(\ref{eq:phi_interior}) carries an explicit
  $e^{-\phi}$ prefactor on its matter-source side and
  is \emph{not} strictly shift-invariant in matter; shift
  invariance of the full coupled system is recovered only
  at leading order in $|\phi|$ (Sec.~\ref{sec:shift_coupled}).

\item The kinematic channel selects the value
  \begin{equation}
    s = \frac{1}{2}\quad\text{(leading order, for all
    bodies)}
  \end{equation}
  via the bare-mass Komar integrand
  $\rho_0\,e^{-\phi} = \rho_0\,e^{\phi/2}\,e^{-3\phi/2}$,
  which combines the universal conformal redshift
  $e^{\phi/2}$ with the proper-volume factor $e^{-3\phi/2}$
  (Sec.~\ref{sec:s_kinematic}); the redshift factor alone
  does not produce $s=1/2$, since the formal
  definition~(\ref{eq:s_formal}) acts on $M_{\mathrm{Komar}}$,
  not on the test-particle Killing-energy
  $m_{\mathrm{eff}}$.  At leading order in $|\phi|$, the
  excess sensitivity satisfies
  \begin{equation}
    s_{\mathrm{excess}} = 0 + O(|\phi|)\,.
  \end{equation}

\item A naive perturbation analysis that holds $\rho$
  fixed while varying $\phi$ yields a screening term
  $\mu^2\,\delta\phi$ of the wrong structural type.  The
  correct coupled linearisation (Sec.~\ref{sec:correction})
  shows that this term is \emph{not a solution} of the
  coupled system in matter, but it also shows that a
  constant-$\delta\phi$ rigid-shift ansatz with
  $\delta\rho = \delta p = 0$ is \emph{not} a solution
  inside the star either.  Whether a modified ansatz
  with self-consistent $(\delta\rho,\delta p)$ solves the
  coupled system non-perturbatively at strong field is
  an \textbf{open task}.

\item The position-dependent effective mass
  $m_{\mathrm{eff}} = m_0\,e^{\phi/2}$ is fully retained;
  it is one of the two factors (alongside the
  proper-volume element $e^{-3\phi/2}$) entering the
  bare-mass Komar integrand $\rho_0\,e^{-\phi}$ that
  produces the leading-order universality of $s=1/2$
  (Sec.~\ref{sec:s_kinematic}).  The theory does not
  require abandoning the
  variable mass.

\item Vacuum bi-conformal compact objects carry
  $s_{\mathrm{excess}} = 0$ exactly, as derived from
  the vacuum exterior harmonic monopole profile (main theory
  Sec.~7) \emph{combined with} the rigid-shift response
  of the bare-mass Komar integral
  (Sec.~\ref{sec:nohair}, item~1).  This is the only
  regime in which the exact statement is presently
  proven.

\item The theory is compatible with all current
  pulsar-timing tests at leading order:
  \begin{itemize}
  \item J0337+1715 (Nordtvedt): $\Delta a/a \simeq 0$
    at leading order vs.\ conservative 2018
    Archibald bound $|\Delta a/a| < 2.6\times10^{-6}$
    (tightened to model-dependent
    $|\Delta|\lesssim 1.5$--$2.3\times10^{-6}$ by the
    2025 Voisin et~al.\ re-analysis~\cite{Voisin2025})
    --- consistent if the open task closes with
    $s_{\mathrm{excess}}\to 0$ to the required precision.
  \item J1738+0333 (orbital decay):
    $\dot{P}_b^{(\mathrm{dipole})} \simeq 0$ at leading
    order vs.\ illustrative bound
    $|s_{\mathrm{excess}}| \lesssim 0.05$
    (pending the ISPG-specific $\kappa_D$ normalization;
    see Sec.~\ref{sec:bounds} Compatibility-table
    footnote) --- consistent.
  \item LIGO (BH--BH binaries):
    no dipole radiation --- consistent, rigorously
    (vacuum no-hair).
  \end{itemize}

\item The structural reason for the leading-order
  universality is unique to the bi-conformal framework:
  constant $G_4$ (no Brans--Dicke running) and minimal
  coupling ensure that $d\phi/dr$, not $\phi$, governs
  the interior hydrostatic equilibrium.  This is what
  distinguishes the theory from Brans--Dicke, DEF, and
  TeVeS, all of which have $s_{\mathrm{excess}} = O(1)$
  compactness-dependent effects.

\item \textbf{Quantum-paper conjecture and open task.}
  The companion quantum paper~\cite{Lotishvili2026Q}
  conjectures that the leading-order classical shift
  invariance of the TOV sector is promoted to an exact
  Ward identity of the full quantised coupled system,
  which would make $s = 1/2$ and $s_{\mathrm{excess}} = 0$
  exact at all loop orders.  The present appendix does
  \emph{not} prove this: the classical coupled system
  as written has only leading-order shift invariance
  in matter.  A rigorous derivation of an exact rigid-shift
  symmetry (classical or quantum) for the matter sector
  is the principal open task of the sensitivity programme.
\end{enumerate}
\end{document}
